% =====================================================================
\chapter{前日準備}
% =====================================================================

前日にやることは，\textbf{「作ったデータが正しいことを実物で確かめる」}ことと，
\textbf{「壊れたときの逃げ道を用意する」}ことの 2 つです。
新しくデータを作る作業は，前日に残さない形にしておきます。

% ---------------------------------------------------------------------
\Step{試走・ポ確でコースデータを検証する}{\tALL\ \tSI\ \tEMIT}

\begin{Goal}
実際にコントロールを回ったカードを読み，
\textbf{全コントロール番号が，コースデータどおりの順番で記録されている}ことを確認する。
\end{Goal}

事前準備のなかで，ここがいちばん効く作業です。
ここを通れば，コースデータの誤り・ユニットの設定ミス・設置場所の入れ違い・
電池切れが，まとめて発見できます。

\subsection*{手順}

\begin{enumerate}
\item 試走者・ポ確者に\textbf{専用のカード}を渡します（競技用のレンタルカードとは分けます）。
      \tEMIT\ EMIT の場合は，出走前に\textbf{スタートユニットでアクティベート}してもらう必要があるので，
      スタートユニットも渡します。
\item \tSI\ 併せて\textbf{ピンパンチ用の紙}（地図のリザーブ欄に相当するもの）を渡し，
      各コントロールでピンパンチも押してもらいます。
      これが\textbf{ピンパンチ対応表}（Step 18）の原本になります。
      \tEMIT\ EMIT ではこの作業は不要です（ユニット自体がラベルに穴を開けるため）。
\item 帰ってきたカードを Mulka2 で読みます。
      メインウィンドウ →【EMIT/SI】（バージョンによっては【ツール】）→
      \tEMIT\ 【E カード読み取り】／\tSI\ 【SI カード読み取り】→
      モード\textbf{【試走カード読み取り】}を選び，読み取り機に接続します。\ver
      COM 番号が分からないときは，Windows の\textbf{デバイスマネージャ}でも確認できます。
\item 読み取り結果を，\textbf{回ってきた本人と一緒に}確認します。
      \begin{itemize}
      \item コースが自動認識されているか
      \item 記録された番号の\textbf{順番}が，本人の記憶どおりか
      \item 抜けている番号がないか
      \item \tEMIT\ 電池警告が出ていないか
      \end{itemize}
\item \tSI\ \tSIAC\ タッチフリーを使う場合は，起動を\textbf{差し込みパンチ}で行い，
      そのうえで\textbf{全コントロールをタッチフリーでもパンチ}して反応を確認します。
\item 全コントロールが確認できたら，\textbf{競技責任者に「競技開始可能」と伝えます}。
\end{enumerate}

\begin{Note}[試走で「フィニッシュ不通過」と出ても異常とは限らない]
試走の時点で\textbf{パンチングフィニッシュのユニットをまだ設置していない}場合，
試走者はフィニッシュをパンチできません。
そのため読み取ると\textbf{フィニッシュ不通過}と表示されますが，これは想定どおりです。
フィニッシュ機材の確認は，設置後に別途行います。

同じく，コースデータと試走内容が違えばミスパンチの表示が出ます。
このとき\textbf{「コースデータが間違っているのか，試走の回り方が違ったのか」}を
切り分けてから直します。どちらかを決めずに直すと，正しいほうを壊します。
\end{Note}

\begin{Note}[Mulka2 を使わずに確認する方法]
試走・ポ確の責任者が計センから離れている場合，
\tEMIT\ \textbf{MTR ＋プリンタ}，\tSI\ \textbf{メインステーション＋プリントセット}でも
その場で結果を確認できます。
MTR は持ち運びやすく操作も簡単なので，
設置担当が 1 台持っていれば，計センに戻らずに確認が完結します。
\end{Note}

\begin{Note}[「本人と一緒に」確認するのはなぜか]
設置者がコントロール A と B を取り違えて置いた場合，
カードのデータは「順番どおりでない番号の列」になります。
しかし\textbf{画面だけ見ていると「そういうコースなのかな」と流してしまいます}。
回ってきた本人に\textbf{記憶どおりの順番}を口で言ってもらい，
それと画面が一致するかを照合すると，設置ミスをここで止められます。
食い違ったら，本人の記憶違いでない限り\textbf{設置ミス}です。
\end{Note}

\subsection*{ユニット／ステーションを交換することになったら}

\begin{enumerate}
\item 予備の機材に，\textbf{故障した機材と同じ番号}を設定します（\tSI\ Config+ で設定）。
      これができれば，コースデータも地図も変更不要です。これが第一選択です。
\item 同じ番号にできない場合は，\textbf{別番号の機材を使い，Mulka2 側を変更}します。
      メインウィンドウ →【入力/編集】→【コース設定変更】から，
      該当番号を\textbf{一括で}置き換えます。
\item 地図が印刷済みの場合は，\textbf{新しい番号の上に古い番号のシールを貼って偽装}します。
      競技者は地図どおりの番号を探すためです。
\item \textbf{公式掲示板・注意事項で「地図の表記と実際の番号が異なるコントロールがある」ことを周知}します。
\end{enumerate}

\begin{Done}
全コースについて，試走カードの読み取り結果とコースデータが一致した。
競技責任者に報告済み。
\end{Done}

\begin{Fail}
\textbf{起きること：}試走はしたが，\textbf{フィニッシュのパンチを確認していない}。\\
\textbf{対処：}フィニッシュのユニット／ステーションが壊れていると，
\textbf{全員のフィニッシュ時刻が取れなくなります}。
試走カードには必ずフィニッシュのパンチも入れ，読み取り画面でフィニッシュ時刻が
出ていることを確認します。出ていない場合は，
(1) 機材の故障，(2) Mulka2 側でパンチングフィニッシュになっていない（「記録を計算できません」になる），
(3) フィニッシュのユニット番号がコースデータと違う（DNF になる）
のいずれかです。3 つとも Mulka2 側で直せます。
\end{Fail}

% ---------------------------------------------------------------------
\Step{バックアップ計時を準備する}{\tALL}

\begin{Goal}
「カードが読めなかった人」の記録を，計セン以外の手段で復元できる体制を作る。
\end{Goal}

\begin{Note}[バックアップ計時は「使うもの」です]
カードが読めない事態は，参加者 100 人あたり 1 人程度の割合で起きます。
\textbf{100 人規模の大会なら，1 回は使う}計算になります。
使わないだろうと考えて省くと，そのときに記録が出せません。
\end{Note}

\subsection*{(1) 通過の証明（何番を通ったか）}

\begin{center}
\small
\begin{tabular}{@{}P{20mm}P{58mm}P{62mm}@{}}
\toprule
 & \hd{\tEMIT} & \hd{\tSI\ \tSIAC} \\
\midrule
手段 & \textbf{バックアップラベル（BUL）} & \textbf{地図リザーブ欄へのピンパンチ} \\
仕組み & \textbf{ユニット自体}が，カードに挟んだラベルに穴を開ける。
  穴のパターンはユニット固有
  & コントロールに\textbf{別途ピンパンチを吊るし}，競技者が地図の欄に穴を開ける \\
当日の準備 & \textbf{ラベルがカードに入っていることを確認するだけ}
  & \textbf{ピンパンチ対応表を作る}（次項） \\
読み方 & Mulka2 のペナレポートに「正解ならどう開くか」が印刷されるので，
  それと実物を照合する（Step 31）
  & 対応表と穴の位置を照合する \\
\bottomrule
\end{tabular}
\end{center}

\begin{Note}[\tEMIT\ EMIT では対応表を当日作る必要がない]
EMIT は\textbf{ユニットそのものがバックアップラベルに穴を開けます}。
どのユニットがどのパターンを開けるかは\textbf{機材に固有}なので，
「どれをどこに吊るしたか」を当日記録する作業がありません。

照合は，同じコースの\textbf{完走者}の【競技者情報】→【カード詳細・対応】を開けば，
そのコースの\textbf{正解の穴の並び}が表示されます。
それと，読めなかった人の実物のラベルを見比べます。
\end{Note}

\paragraph{\tSI\ ピンパンチ対応表の作り方。}
\textbf{この作業は SI のときだけ必要です。}
SI にはバックアップラベルがないので，コントロールに\textbf{ピンパンチを別に吊るします}。
ピンパンチの穴のパターンは機材ごとに固有なので，
\textbf{どのコントロールにどのパターンを吊るしたか}を記録しないと，あとで照合できません。

\begin{enumerate}
\item 設置者に，各コントロールでピンパンチを紙に押してもらう（Step 17 と同時に行う）。
\item 紙に\textbf{コントロール番号を書き込んで}回収する。
\item これをコピーして\textbf{計セン用の対応表}にする。当日朝の作成で構いません。
\end{enumerate}

\begin{Fig}
\begin{tikzpicture}[x=1mm,y=1mm]
  \node[scr,minimum width=140mm,minimum height=40mm,anchor=north west] (w) at (0,0){};
  \node[anchor=north west,font=\sffamily\bfseries\small] at (3,-3) {ピンパンチ対応表（SI のみ。当日朝に設置者から回収して作成）};
  \foreach \i/\n in {0/31,1/32,2/33,3/34,4/35,5/36}{
    \draw[boxln,line width=0.4pt] (6+\i*22,-10) rectangle (24+\i*22,-30);
    \node[anchor=north,font=\sffamily\bfseries\small] at (15+\i*22,-32.5) {\n};
  }
  % random-ish pin patterns
  \foreach \i/\pa/\pb/\pc/\pd in {0/2/5/9/13, 1/1/6/10/15, 2/3/4/11/14, 3/2/7/8/16, 4/5/6/12/13, 5/1/8/11/15}{
    \foreach \p in {\pa,\pb,\pc,\pd}{
      \pgfmathsetmacro{\cx}{9 + \i*22 + mod(\p-1,4)*4}
      \pgfmathsetmacro{\cy}{-13 - floor((\p-1)/4)*4}
      \fill[boxln] (\cx,\cy) circle (0.9);
    }
  }
  \node[anchor=north west,font=\sffamily\tiny] at (3,-35) {※ 実際のパターンは機材ごとに固有。上図は配置イメージ};
\end{tikzpicture}
\figcap{ピンパンチ対応表（\tSI\ SI のみ。穴のパターンとコントロール番号の対応）}
\end{Fig}

\subsection*{(2) 時刻の記録（何時にフィニッシュしたか）}

通過の証明ができても，\textbf{フィニッシュ時刻}が分からなければ記録は出せません。
フィニッシュパートと分担を決めます。方法は 3 つあり，どれか 1 つ以上を動かします。

\begin{center}
\small
\begin{tabular}{@{}P{34mm}P{50mm}P{56mm}@{}}
\toprule
\hd{方法} & \hd{利点} & \hd{欠点・注意} \\
\midrule
プリンタ付きストップウォッチ ＋ 着順シール
  & その場でレシートを渡せる。すぐ照合できる
  & 人手が 2 名必要（押す人・シールを貼る人）。紙切れ・紙詰まりに注意，とくに雨 \\
定点カメラ（動画）
  & 人手 0 で運用できる
  & 閉鎖後にカメラを回収しないと確認できない。ゼッケン着用がほぼ必須 \\
スマートフォンで撮影時刻を記録
  & 1 名で足りる。競技中に確認できる
  & 撮影時刻を記録するアプリが必要。端末の時計を合わせておく \\
\bottomrule
\end{tabular}
\end{center}

\begin{Note}[最低限これだけは]
凝った仕組みが用意できないときは，フィニッシュ係が紙に
\textbf{「フィニッシュ時刻」と「カード番号」を書き取る}だけで十分機能します。
計センから「30 分ほど前にフィニッシュした，カード番号 8100123 の人の時刻がほしい」
と聞かれたときに答えられれば，目的は達しています。
\end{Note}

\begin{Done}
\tSI\ ピンパンチ対応表ができている（\tEMIT\ EMIT はラベルの有無を確認した）。
フィニッシュでの予備計時の担当者と方法が決まっている。
\end{Done}

% ---------------------------------------------------------------------
\Step{データをバックアップして共有する}{\tALL}

\begin{Goal}
チーフの PC が壊れても，別の PC で大会が回せる状態にする。
\end{Goal}

\begin{enumerate}
\item Mulka2 のデータフォルダから，\textbf{該当イベントのフォルダを丸ごと zip}にします。
\item 計センメンバー全員に配ります（メッセージアプリ・USB メモリ・クラウドストレージ）。
\item 受け取った側は，自分の PC の Mulka2 データフォルダに\textbf{解凍して置く}だけで，
      同じイベントが開けます。
\item \textbf{各 PC で実際に開いて，メインウィンドウが正常に起動することを確認}します。
\end{enumerate}

\begin{Fig}
\begin{tikzpicture}[node distance=6mm]
  \node[blkd,minimum width=30mm,minimum height=11mm] (c) {チーフ PC\\\scriptsize イベントデータ原本};
  \node[blkg,minimum width=26mm,right=16mm of c,yshift=11mm] (a) {メンバー A の PC};
  \node[blkg,minimum width=26mm,right=16mm of c] (b) {メンバー B の PC};
  \node[blkg,minimum width=26mm,right=16mm of c,yshift=-11mm] (d) {予備 PC ／ USB メモリ};
  \draw[ar] (c)--node[lbl]{zip}(a);
  \draw[ar] (c)--node[lbl]{zip}(b);
  \draw[ar] (c)--node[lbl]{zip}(d);
  \node[blk,below=10mm of c,minimum width=52mm,align=left,font=\sffamily\scriptsize]
    (n) {\textbf{ファイル名のルール}\\
         ・改訂日時を入れる（例：\texttt{20260510\_1930\_v3}）\\
         ・古いものは \texttt{old\textbackslash{}} に移す\\
         ・「最新はどれか」を口頭で確認しない（名前で分かるようにする）};
\end{tikzpicture}
\figcap{イベントデータの共有}
\end{Fig}

\begin{Done}
計センメンバー全員の PC で，同じイベントデータが開ける。
\end{Done}

\begin{Fail}
\textbf{起きること：}古いデータを誤って使う。これはよく起きます。\\
\textbf{対処：}\textbf{当日朝，全 PC で「エントリー人数」を突き合わせます}。
人数が違えば版が違います。1 分で終わる確認です。
\end{Fail}

% ---------------------------------------------------------------------
\Step{印刷物をそろえる}{\tALL\ \tRANK}

\begin{Goal}
当日に紙で必要になるものを全部刷り，配布先ごとに束ねる。
\end{Goal}

\begin{center}
\small
\begin{tabular}{@{}P{42mm}P{18mm}P{84mm}@{}}
\toprule
\hd{印刷物} & \hd{配布先} & \hd{備考} \\
\midrule
救護用スタートリスト & 救護・パトロール & 全員通し・ゼッケン番号順（Step 13 (a)） \\
スタート用スタートリスト & スタート & レーン別・時刻順（Step 13 (b)）。レーン数 $\times$ 2 部 \\
フィニッシュ用リスト & フィニッシュ & ゼッケン番号順・帰還チェック欄付き（Step 13 (c)） \\
公開用スタートリスト & 会場掲示・スタート地区掲示 & 大きめに印刷 \\
レンタルカードリスト & 計セン・受付 & カード番号順。回収チェック欄付き \\
当日申込用の空き枠一覧 & 受付・計セン & ゼッケン番号と時刻だけ埋まった状態 \\
運営用全ポ図 & 計セン & \textbf{ペナチェックで必須}。3 セット程度 \\
各コース図 & 計セン & 同上 \\
ピンパンチ対応表 & 計セン & 当日朝に作る場合は用紙だけ準備 \\
カード変更届の用紙 & 受付 & 「ゼッケン番号・氏名・新しいカード番号」の 3 項目でよい \\
\bottomrule
\end{tabular}
\end{center}

\begin{Note}[全ポ図とコース図は紙で持つ]
ペナチェック（Step 30）では，
「あなたはここを通らずにここへ行っています」と\textbf{地図の上で指し示して}説明します。
画面を見せるより圧倒的に速く，納得も得やすいです。
\end{Note}

\begin{Done}
上表の全項目が印刷され，配布先ごとに束ねてある。
各紙に版数と発行時刻が入っている。
\end{Done}

% ---------------------------------------------------------------------
\Step{積み込みと最終確認}{\tALL}

\begin{Goal}
翌朝そのまま設営を始められる状態にする。
\end{Goal}

\begin{enumerate}
\item 機材リスト（Step 2）で\textbf{実物を数えながら}積み込む。
\item \textbf{モバイル Wi-Fi の通信を実際に試す}。
      Mulka2 のクラウドに\textbf{テスト用サーバ}で接続してみるのが確実です。
      本番用サーバに接続してしまうと成績が公開される場合があるので，
      準備段階では\textbf{テスト用サーバ}を使います。\ver
\item PC を全台充電する。電源が取れない会場なら，モバイルバッテリーも充電する。
\item プリンタで\textbf{1 枚試し刷り}する。
\item 計セン連絡用のグループ（会場・受付／スタート／フィニッシュ／救護／計セン）を作る。
      \textbf{同時に，各パートチーフの電話番号を控えます}。
\end{enumerate}

\begin{Note}[連絡手段は「音声通話」を残しておく]
テレインではデータ通信が不安定になりがちです。
メッセージアプリの無料通話はデータ通信を使うため，聞き取れないことがあります。
\textbf{音声通話は別の回線}なので，データが不安定でも通じる可能性があります。
そのため，\textbf{電話番号の交換}をしておきます。
運営スタッフの名札の裏に電話番号を書いておく，という方法もあります。
\end{Note}

\begin{Done}
機材リストの全項目が積み込まれ，Wi-Fi・PC・プリンタの動作を確認した。
各パートチーフの電話番号を控えた。
\end{Done}

\begin{Fail}
\textbf{起きること：}PC の持ち主が設置や試走に行ってしまい，
\textbf{ログインパスワードが分からない}。\\
\textbf{対処：}前日のうちに，各 PC のパスワードを紙に書いて
チーフが預かります。当日朝に「持ち主が山にいて PC が開けない」となると，
その PC は 1 日使えません。
\end{Fail}
